\section{Conclusion}

This paper demonstrated the feasibility of opportunistic Integrated Sensing and Communications (ISAC) for passenger movement classification in public transport using commodity Wi-Fi infrastructure. By passively capturing standard \ac{RSSI} sequences without dedicated radar hardware or firmware modifications, self-supervised \ac{JEPA} world models learn environment-invariant representations of passenger mobility. On the extended 3,511-sample dataset, TS-JEPA achieves a single-model \ac{MCC} of 0.9194, while a 16-member calibrated stacked ensemble reaches 0.9232. On the earlier collection, a compact three-member Pareto ensemble attains 0.8904 \ac{MCC} with only 1.26~M parameters, substantially surpassing the prior classical Gaussian Process baseline of 0.756.

A central insight is the stark contrast in cross-environment transferability between self-supervised and supervised sequence modeling. While self-supervised world models exhibit positive transfer across collection environments with \ac{MCC} gains between +0.029 and +0.042, supervised regularized baselines suffer representational collapse. Furthermore, systematic permutation importance reveals that a compact five-channel feature subset retains 95\% of total discriminative power, enabling sub-millisecond front-end computation and real-time execution directly on resource-constrained edge access points.

These findings confirm that commodity Wi-Fi access points can effectively serve as dual-functional nodes, delivering communications while providing non-intrusive, privacy-preserving sensing layers for intelligent transportation systems. Future research will explore cooperative multi-access-point fusion across vehicle doors for spatial diversity, long-horizon temporal modeling across full transit routes, and online continual adaptation algorithms to compensate for seasonal radio propagation drift. Open-source models and reproducibility scripts are available as supplementary material \cite{githubrepo}.
